\documentclass[10pt,a4paper]{amsart}
\usepackage{amsmath,amssymb,amsthm,graphicx}
% For commutative diagrams we use Xy-pic
\usepackage[all]{xy}

% We use 2cell for 2-commutative diagrams.
\xyoption{2cell}
% Theorem environments.
%
\theoremstyle{plain}
\newtheorem{theorem}[subsection]{Theorem}
\newtheorem{proposition}[subsection]{Proposition}
\newtheorem{lemma}[subsection]{Lemma}

\theoremstyle{definition}
\newtheorem{definition}[subsection]{Definition}
\newtheorem{example}[subsection]{Example}
\newtheorem{exercise}[subsection]{Exercise}
\newtheorem{situation}[subsection]{Situation}

\theoremstyle{remark}
\newtheorem{remark}[subsection]{Remark}
\newtheorem{remarks}[subsection]{Remarks}

\numberwithin{equation}{subsection}

% Macros
%
\def\lim{\mathop{\mathrm{lim}}\nolimits}
\def\colim{\mathop{\mathrm{colim}}\nolimits}
\def\Spec{\mathop{\mathrm{Spec}}}
\def\Hom{\mathop{\mathrm{Hom}}\nolimits}
\def\Ext{\mathop{\mathrm{Ext}}\nolimits}
\def\SheafHom{\mathop{\mathcal{H}\!\mathit{om}}\nolimits}
\def\SheafExt{\mathop{\mathcal{E}\!\mathit{xt}}\nolimits}
\def\Sch{\mathit{Sch}}
\def\Mor{\mathop{\mathrm{Mor}}\nolimits}
\def\Ob{\mathop{\mathrm{Ob}}\nolimits}
\def\Sh{\mathop{\mathit{Sh}}\nolimits}
\def\NL{\mathop{N\!L}\nolimits}
\def\CH{\mathop{\mathrm{CH}}\nolimits}
\def\proetale{{pro\text{-}\acute{e}tale}}
\def\etale{{\acute{e}tale}}
\def\QCoh{\mathit{QCoh}}
\def\Ker{\mathop{\mathrm{Ker}}}
\def\Im{\mathop{\mathrm{Im}}}
\def\Coker{\mathop{\mathrm{Coker}}}
\def\Coim{\mathop{\mathrm{Coim}}}

% Boxtimes
%
\DeclareMathSymbol{\boxtimes}{\mathbin}{AMSa}{"02}

%
% Macros for moduli stacks/spaces
%
\def\QCohstack{\mathcal{QC}\!\mathit{oh}}
\def\Cohstack{\mathcal{C}\!\mathit{oh}}
\def\Spacesstack{\mathcal{S}\!\mathit{paces}}
\def\Quotfunctor{\mathrm{Quot}}
\def\Hilbfunctor{\mathrm{Hilb}}
\def\Curvesstack{\mathcal{C}\!\mathit{urves}}
\def\Polarizedstack{\mathcal{P}\!\mathit{olarized}}
\def\Complexesstack{\mathcal{C}\!\mathit{omplexes}}
% \Pic is the operator that assigns to X its picard group, usage \Pic(X)
% \Picardstack_{X/B} denotes the Picard stack of X over B
% \Picardfunctor_{X/B} denotes the Picard functor of X over B
\def\Pic{\mathop{\mathrm{Pic}}\nolimits}
\def\Picardstack{\mathcal{P}\!\mathit{ic}}
\def\Picardfunctor{\mathrm{Pic}}
\def\Deformationcategory{\mathcal{D}\!\mathit{ef}}
\newenvironment{reference}{}{}
\newenvironment{slogan}{}{}
\newenvironment{history}{}{}
\usepackage{cleveref}
\crefname{originalsection}{original section}{original sections}
\Crefname{originalsection}{Original section}{Original sections}
\makeatletter
\newlabel{section-semistable-reduction-genus-zero}{{0CDK}{source}{}{sourcepointer-1}{}}
\newlabel{section-semistable-reduction-genus-zero@cref}{{[originalsection][0][]0CDK}{[1][0][]source}{}{}{}}
\makeatother
\crefname{originalsection}{original section}{original sections}
\Crefname{originalsection}{Original section}{Original sections}
\makeatletter
\newlabel{section-semistable-reduction-genus-one}{{0CEG}{source}{}{sourcepointer-2}{}}
\newlabel{section-semistable-reduction-genus-one@cref}{{[originalsection][0][]0CEG}{[1][0][]source}{}{}{}}
\makeatother
\crefname{originalsection}{original section}{original sections}
\Crefname{originalsection}{Original section}{Original sections}
\makeatletter
\newlabel{section-semistable-reduction-genus-at-least-two}{{0CEI}{source}{}{sourcepointer-3}{}}
\newlabel{section-semistable-reduction-genus-at-least-two@cref}{{[originalsection][0][]0CEI}{[1][0][]source}{}{}{}}
\makeatother
\crefname{originallemma}{original lemma}{original lemmas}
\Crefname{originallemma}{Original lemma}{Original lemmas}
\makeatletter
\newlabel{curves-lemma-torsion-picard-becomes-visible}{{0CDU}{source}{}{sourcepointer-4}{}}
\newlabel{curves-lemma-torsion-picard-becomes-visible@cref}{{[originallemma][0][]0CDU}{[1][0][]source}{}{}{}}
\makeatother
\crefname{originalequation}{original equation}{original equations}
\Crefname{originalequation}{Original equation}{Original equations}
\makeatletter
\newlabel{equation-bound}{{0CEJ}{source}{}{sourcepointer-5}{}}
\newlabel{equation-bound@cref}{{[originalequation][0][]0CEJ}{[1][0][]source}{}{}{}}
\makeatother
\title{Semistable reduction for a smooth projective curve after a finite separable extension of the base}
\author{The Stacks Project Authors}
\date{stacks-project:models.tex:local-snapshot-20261012; GNU FDL 1.2 or later}
\begin{document}
\maketitle
\noindent\textit{Editorial note.} The counted unit is the Stacks Project Theorem theorem-semistable-reduction at source lines 6856-6885 together with its proof at source lines 6887-6900, a straight adjacent pairing in the chapter file. The chapter is the upstream file models.tex; the pool directory carries the longer name stacks-models and is not used as the chapter name anywhere in this card.\par\medskip
\noindent\textit{Editorial note.} The statement opens with the author's reference environment, which the retained wrapper renders as a visible no-op and which quotes the cited corollary on the irreducibility of the space of curves of a given genus. That citation is the only one in this card, and its entry is supplied as a task-directory fragment holding it exactly as a bibtex run on the host had formatted it out of the Stacks Project's shared bibliography; the bibliography backend is therefore none.\par\medskip
\noindent\textit{Editorial note.} This card has no context span: the statement and proof are self-contained apart from the five pointers, and the notions they use are recorded in the standard prerequisites.\par\medskip
\noindent\textit{Editorial note.} The retained proof names three sections of this chapter and two results of the Curves chapter. All are declared as non-load-bearing navigation pointers with their stable Stacks tags looked up in the shared tag index, and none of their statements is reproduced inside the span.\par\medskip
\noindent\textit{Editorial note.} The retained proof span starts at line 6887, the line of the opening marker read from the source, and contains exactly one proof environment.\par\medskip
\noindent\textit{Allowed standard prerequisites (editorial).}\begin{itemize}
\item Models of a curve, semistable reduction and the three genus cases. A model of a smooth projective curve over the fraction field of a discrete valuation ring is a flat proper scheme over the ring with that curve as generic fibre, and its reduction is semistable when the special fibre is nodal; after a finite separable extension of the base one may pass to the integral closure and consider models over its localisations. The chapter treats the three genus cases in three separate sections, proving that a curve of genus zero becomes the projective line after a quadratic separable extension, that a curve of genus one acquires semistable reduction whose special fibre is controlled by the minimal model, and that a curve of genus at least two acquires semistable reduction after a finite separable extension, the last case using the fact that torsion in the Picard group of the model becomes visible after such an extension. The author's own sections and the quoted result of the Curves chapter are the allowed inputs.. The counted theorem is the main result of the chapter: it is the semistable reduction statement that the chapter's preceding sections are built to prove, and its proof is the assembly of those sections.
\end{itemize}\medskip
\ifdefined\hypertarget\hypertarget{sourcepointer-1}{}\fi
\noindent\textit{Source reference.} Original section 0CDK: focused-seeds/AGP/sources/stacks-snapshot-20261011/tags.txt, source lines 15545--15545. Navigation pointer to the chapter's section on the genus-zero case, from which the first enumerated assertion is taken; the cited material is an allowed input and is not reproduced here. The number is the stable Stacks tag of the target, looked up in the shared tag index.\par\medskip
\ifdefined\hypertarget\hypertarget{sourcepointer-2}{}\fi
\noindent\textit{Source reference.} Original section 0CEG: focused-seeds/AGP/sources/stacks-snapshot-20261011/tags.txt, source lines 15576--15576. Navigation pointer to the chapter's section on the genus-one case, from which the second enumerated assertion is taken; the cited material is an allowed input and is not reproduced here. The number is the stable Stacks tag of the target, looked up in the shared tag index.\par\medskip
\ifdefined\hypertarget\hypertarget{sourcepointer-3}{}\fi
\noindent\textit{Source reference.} Original section 0CEI: focused-seeds/AGP/sources/stacks-snapshot-20261011/tags.txt, source lines 15578--15578. Navigation pointer to the chapter's section on the remaining case, from which the third enumerated assertion is taken; the cited material is an allowed input and is not reproduced here. The number is the stable Stacks tag of the target, looked up in the shared tag index.\par\medskip
\ifdefined\hypertarget\hypertarget{sourcepointer-4}{}\fi
\noindent\textit{Source reference.} Original lemma 0CDU: focused-seeds/AGP/sources/stacks-snapshot-20261011/tags.txt, source lines 15554--15554. Allowed input of the Curves chapter quoted by the retained proof for the visibility of torsion in the Picard group after a finite separable extension of the base, which is the key step of the last case; the pointer is registered so the reference resolves and the cited statement is not reproduced inside the span. The number is the stable Stacks tag of the target, looked up in the shared tag index.\par\medskip
\ifdefined\hypertarget\hypertarget{sourcepointer-5}{}\fi
\noindent\textit{Source reference.} Original equation 0CEJ: focused-seeds/AGP/sources/stacks-snapshot-20261011/tags.txt, source lines 15579--15579. Navigation pointer to a labelled displayed bound of the chapter that the retained proof refers to when it recalls the numerical estimate used in the last case; the cited labelled equation is an allowed input and its text is not reproduced here. The number is the stable Stacks tag of the target, looked up in the shared tag index.\par\medskip
\begin{theorem}
\label{theorem-semistable-reduction}
\begin{reference}
\cite[Corollary 2.7]{DM}
\end{reference}
Let $R$ be a discrete valuation ring with fraction field $K$. Let $C$ be a
smooth projective curve over $K$ with $H^0(C, \mathcal{O}_C) = K$.
Then there exists an extension of discrete valuation rings
$R \subset R'$ which induces a finite separable extension of
fraction fields $K'/K$ such that $C_{K'}$ has semistable reduction.
More precisely, we have the following
\begin{enumerate}
\item If the genus of $C$ is zero, then there exists a degree $2$
separable extension $K'/K$ such that $C_{K'} \cong \mathbf{P}^1_{K'}$
and hence $C_{K'}$ is isomorphic to the generic fibre of the
smooth projective scheme $\mathbf{P}^1_{R'}$ over the integral closure
$R'$ of $R$ in $K'$.
\item If the genus of $C$ is one, then there exists a finite separable
extension $K'/K$ such that $C_{K'}$ has semistable reduction
over $R'_\mathfrak m$ for every maximal ideal $\mathfrak m$
of the integral closure $R'$ of $R$ in $K'$. Moreover, the special
fibre of the (unique) minimal model of $C_{K'}$ over $R'_\mathfrak m$
is either a smooth genus one curve or a cycle of rational curves.
\item If the genus $g$ of $C$ is greater than one, then there exists a
finite separable extension $K'/K$ of degree at most
$B_g$ (\ref{equation-bound}) such that $C_{K'}$ has semistable reduction
over $R'_\mathfrak m$ for every maximal ideal $\mathfrak m$
of the integral closure $R'$ of $R$ in $K'$.
\end{enumerate}
\end{theorem}

\begin{proof}
For the case of genus zero, see
Section \ref{section-semistable-reduction-genus-zero}.
For the case of genus one, see
Section \ref{section-semistable-reduction-genus-one}.
For the case of genus greater than one, see
Section \ref{section-semistable-reduction-genus-at-least-two}.
To see that we have a bound on the degree $[K' : K]$
you can use the bound on the degree of the extension needed
to make all $\ell$ or $\ell'$ torsion visible proved in
Algebraic Curves, Lemma \ref{curves-lemma-torsion-picard-becomes-visible}.
(The reason for using $\ell$ and $\ell'$ is that we need to
avoid the characteristic of the residue field $k$.)
\end{proof}

\begin{thebibliography}{10}

\bibitem{DM}
P.~Deligne and D.~Mumford.
\newblock The irreducibility of the space of curves of given genus.
\newblock {\em Publ. Math. IHES}, 36:75--109, 1969.

\end{thebibliography}
\end{document}
